\chapter{Prefazione}

La maggior parte dei corsi di architettura dei calcolatori è costruita dal
basso verso l'alto. Si comincia dal bit, si sale attraverso i sistemi di
numerazione, l'algebra di Boole e le reti logiche, si assembla un processore a
partire dalle porte, e solo nelle ultime settimane si lascia che lo studente
scriva un programma per la macchina che i mesi precedenti hanno descritto.
L'ordine è logicamente impeccabile, e per molti studenti del primo anno è anche
la ragione per cui la materia sembra lontana: per dieci settimane ogni risposta
arriva prima della sua domanda.

Questo libro è la forma scritta di un corso che rovescia l'ordine senza
rinunciare in nulla al rigore. \emph{Prima programmare, poi formalizzare}:
nella prima settimana il lettore installa una piccola console, scrive una
dozzina di righe di assembly \riscv e vede un rettangolo muoversi su uno
schermo. Da quel momento ogni argomento formale è introdotto da qualcosa che il
programma del lettore ha già fatto. Il complemento a due arriva dopo che una
posizione è uscita dal bordo sinistro dello schermo e il debugger ha mostrato
\hex{FFFFFFFF}. L'algebra di Boole arriva dopo che una maschera è servita a
chiedere se un pulsante è premuto. Il flip-flop arriva dopo che un contatore è
sopravvissuto da un frame al successivo, e il ciclo dell'istruzione dopo che il
lettore ha premuto \key{Step} qualche centinaio di volte e ha visto il program
counter avanzare di quattro.

Dal programma classico non viene tolto nulla. Il lettore troverà la notazione
posizionale e i formati IEEE~754, assiomi e teoremi dell'algebra di Boole con
le loro dimostrazioni, forme canoniche e mappe di Karnaugh, il transistor come
interruttore, sommatori, multiplexer e decodificatori, latch e flip-flop,
macchine a stati finiti, registri e bus, gerarchie di memoria e memorie
permanenti, il datapath e l'unità di controllo, le interruzioni, la scheda
video e le periferiche. Cambia la direzione del ragionamento: ogni capitolo si
apre sulla macchina, la lascia per tutto il tempo che la teoria richiede, e vi
ritorna con una sezione intitolata \emph{Ritorno alla macchina}, in cui le
stesse regole vengono rilette nelle istruzioni di un processore vero.

\section*{La macchina}

Il processore è un core \riscv: l'insieme di istruzioni RV32IMAC, così come lo
realizza l'ESP32\nobreakdash-C6 di Espressif. \riscv è uno standard aperto,
abbastanza piccolo da essere imparato per intero in un semestre e abbastanza
reale da trovarsi in miliardi di dispositivi; sostituisce qui il Motorola~68000
delle edizioni precedenti del corso, che sopravvive solo in qualche confronto
storico.

L'ambiente è \prg, una piattaforma didattica aperta sviluppata all'Università
degli Studi di Napoli «Parthenope». Un programma \prg è una \term{cartuccia}:
un piccolo file di codice macchina nativo che fornisce tre funzioni, una
chiamata una sola volta e due chiamate a ogni frame. La stessa cartuccia gira
su un kit su breadboard, sull'emulatore QEMU e su \prgqt, un'applicazione
desktop il cui debugger integrato mostra i trentadue registri e la memoria
mentre il programma viene eseguito un'istruzione alla volta. Per seguire il
libro non serve una scheda: basta un portatile.

\section*{Il progetto}

Una piattaforma usata solo per gli esercizi è una palestra. \prg è anche
un'officina, e il libro ha un progetto che lo attraversa: \vermone, che si
legge \emph{vermone}, un omaggio a Snake scritto a partire da un file vuoto in
assembly \riscv. Dopo ogni blocco di quattro capitoli, un capitolo tutorial
aggiunge una parte del gioco usando soltanto ciò che quei quattro capitoli
hanno insegnato: prima una testa che attraversa un mondo senza bordi, poi la
sterzata, il cibo e un punteggio, poi un corpo che cresce e le regole che
chiudono una partita, poi cifre, livelli e suono. Un ultimo tutorial confeziona
il gioco finito per il Cartridge Store di \prg. Ogni tutorial stampa ogni riga
del programma, dice quali righe sono nuove e quali sono cambiate, e si ferma a
punti di controllo in cui il programma deve girare e comportarsi in un modo
dichiarato.

\section*{Che cosa significa «verificato» in questo libro}

Un listato che non si assembla insegna la lezione sbagliata. Il codice assembly
di questo libro non è battuto dentro il testo: è letto da file che vengono
assemblati, eseguiti sul QEMU ufficiale e sul runner senza interfaccia di
\prgqt, e confrontati con i risultati attesi da test automatici; le quattro
parti di \vermone sono giocate da un giocatore programmato che scatta anche le
schermate stampate nei capitoli tutorial. Dove un frammento non ha potuto
essere eseguito sulla macchina dell'autore (i frammenti CUDA e Verilog
dell'ultimo capitolo, i passi di installazione per sistemi operativi diversi da
quello usato per scrivere il libro) il testo lo dice.

\section*{Ringraziamenti}

Mariacarla Staffa tiene la metà di laboratorio del corso, in cui gli esercizi
di questo libro si svolgono con matita, carta e tastiera. Ivan Cafiero e Simone
Boscaglia hanno contribuito a \prg da studenti e hanno mostrato che la
piattaforma può essere capita, ed estesa, dalle persone per cui è pensata.
Giuseppe Coviello ha contribuito l'implementazione di riferimento di \vermone:
i quattro programmi che i capitoli tutorial costruiscono riga per riga, e che i
test automatici giocano, sono suoi.
Gli studenti delle edizioni precedenti del corso hanno scritto, con le loro
domande, la maggior parte dei paragrafi che cominciano con «Attenzione».

\begin{flushright}
Raffaele Montella\\
Napoli, ottobre 2026
\end{flushright}
